% =========================================================
% Proof by Contradiction
% =========================================================
\paragraph{Proof by Contradiction.}

To prove $P$, assume $\neg P$ and derive a statement that contradicts
a known truth: a hypothesis, an axiom, or a previously proved theorem.
Since the assumption $\neg P$ led to a false statement, $\neg P$ is
false, and therefore $P$ is true.

\begin{tcolorbox}[colback=propbox, colframe=propborder, arc=2pt,
  left=6pt, right=6pt, top=4pt, bottom=4pt,
  title={\small\textbf{Template: Proof by Contradiction}},
  fonttitle=\small\bfseries]
\textbf{Goal.} $P$.\\[4pt]
Suppose for contradiction that $\neg P$.\\
$\vdots$ \quad [derive consequences; identify the absurdity]\\
This contradicts [name the violated fact].\\
Therefore $P$. \qed
\end{tcolorbox}

\begin{remark}[The logic: why it works]
Contradiction relies on the law of the excluded middle: every statement
is either true or false. To prove $P$, it suffices to show $\neg P$
is false. You show $\neg P$ is false by deriving from it a statement
that is known to be false --- a contradiction.

The contradiction can be anything: $0 = 1$, a set that is both empty
and non-empty, a number that is simultaneously even and odd, an
inequality $x < x$. The only requirement is that the contradicted fact
was genuinely established before the current proof began.

This is the move that students most often get wrong: they think they
have derived a contradiction when they have merely derived something
surprising. A statement being counterintuitive is not a contradiction.
The contradicted fact must be something you can point to:
a hypothesis, an axiom, a definition, or a theorem with a name.
\end{remark}

\begin{remark}[When contradiction is the right tool]
Contradiction is best reserved for claims whose negation is
a strong, productive assumption. The archetypes where contradiction
excels are:

\medskip
\noindent \textbf{Irrationality proofs.} Assuming $\sqrt{2} \in \mathbb{Q}$
gives you a fraction in lowest terms; arithmetic then forces the
numerator and denominator to be simultaneously even, contradicting
the lowest-terms assumption.

\noindent \textbf{Infinitude arguments.} Assuming there are finitely many
primes lets you list them; constructing a number not divisible by any
of them contradicts the assumption that the list was complete.

\noindent \textbf{Nonexistence results.} If you want to show no object
satisfies some property, assume one exists and derive an absurdity.

\noindent \textbf{Uniqueness} (via assume-two-and-compare). Assuming two
distinct objects both satisfy a uniqueness condition produces a
contradiction from the definitions.

\medskip
\noindent
Avoid contradiction as a default. When a direct proof or contrapositive
works, prefer it: the argument is shorter and the logical structure is
more transparent. Contradiction is a sign that the natural direction
of inference does not work, and that you need to reason from an
impossible assumption instead.
\end{remark}

\begin{remark}[Always name the contradiction explicitly]
The closing line of a contradiction proof must name what has been
contradicted. ``This is a contradiction'' with no further specification
is a red flag: it suggests the student either does not know what was
contradicted, or is hoping the reader will not check.

Write: ``This contradicts the assumption that $\gcd(p,q) = 1$.''
Write: ``This contradicts Theorem~X, which asserts $\ldots$''
Write: ``This contradicts $m \in S$ being the least element.''

Naming the violated fact performs two functions: it confirms that the
contradiction is genuine, and it locates the logical break in the
argument so that anyone reading the proof can verify it.
\end{remark}

\begin{example}[$\sqrt{2}$ is irrational]
\textit{Proof.}
Suppose for contradiction that $\sqrt{2} \in \mathbb{Q}$.
Write $\sqrt{2} = p/q$ where $p, q \in \mathbb{Z}$, $q \neq 0$,
and $\gcd(p,q) = 1$ (in lowest terms; this is possible by the
definition of $\mathbb{Q}$).

Squaring: $2q^2 = p^2$.
So $p^2$ is even, which forces $p$ to be even (by the contrapositive
example above). Write $p = 2k$.

Then $2q^2 = (2k)^2 = 4k^2$, so $q^2 = 2k^2$.
So $q^2$ is even, which forces $q$ to be even.

But then both $p$ and $q$ are even, so $\gcd(p,q) \geq 2$.
This contradicts $\gcd(p,q) = 1$.

Therefore $\sqrt{2} \notin \mathbb{Q}$. \qed

\medskip
\noindent
The structure is: assume $\neg P$ (i.e., $\sqrt{2} \in \mathbb{Q}$);
extract information from this assumption (lowest-terms representation);
derive a consequence (both $p$ and $q$ even); identify the contradiction
(conflicts with lowest terms); conclude $P$.
\end{example}
